% =====================================================================
%  Lecture 07 -- LTE Scheduling                           (English edition)
% =====================================================================
\documentclass[11pt,a4paper]{article}
\providecommand\ACNroot{../..}
\input{\ACNroot/templates/lecture-en}

\ACNset{number}{07}
\ACNset{title}{LTE Scheduling}
\ACNset{shorttitle}{LTE Scheduling}
\ACNset{titleother}{زمان‌بندی در LTE}
\ACNset{part}{Part II -- 4G access and core (Lectures 5--13)}
\ACNset{time}{8--10 hours (3.5 h reading, 4--6 h exercises incl.\ simulation)}
\ACNset{level}{Intermediate to advanced}
\ACNset{prereq}{Lectures 5--6; Shannon capacity; Rayleigh fading; basic optimisation}
\ACNset{gen}{4}

\begin{document}
\ACNcover

% ---------------------------------------------------------------------
\section{Learning objectives}
After studying this lecture you should be able to:
\begin{objectives}
  \item \textbf{Describe} the LTE resource grid -- frame, subframe, slot, resource element,
        resource block -- and \textbf{count} the resources of any LTE bandwidth.
  \item \textbf{Explain} downlink and uplink scheduling signalling: DCI, grants, BSR, PHR, CQI.
  \item \textbf{Apply} link adaptation: map CQI to modulation and coding, and \textbf{analyse} outer-loop
        link adaptation (OLLA).
  \item \textbf{Analyse} HARQ with eight stop-and-wait processes and compute its throughput effect.
  \item \textbf{Compare} round robin, maximum C/I and proportional fair scheduling and
        \textbf{derive} why PF maximises the sum of log-throughputs.
  \item \textbf{Evaluate} the throughput--fairness--QoS trade-off with Jain's index and a simulation.
\end{objectives}

% ---------------------------------------------------------------------
\section{Prerequisite knowledge}
You need the LTE protocol stack (Lecture~6), in particular the MAC layer's scheduling and HARQ
roles, the Shannon capacity $C = B\log_2(1+\SNR)$ and the fact that, under Rayleigh fading, the
instantaneous SNR is exponentially distributed around its mean. Basic calculus (derivatives,
Lagrange-style arguments) is used in \cref{sec:pf}.

\paragraph{Key terms} The vocabulary of this lecture with its Persian equivalents:
\begin{keyterms}
Resource element (RE) & \fa{عنصر منبع} & One subcarrier during one OFDM symbol. \\
Resource block (RB) & \fa{بلوک منبع} & 12 subcarriers (180\,kHz) $\times$ one slot. \\
TTI & \fa{بازهٔ زمانی ارسال} & Scheduling interval: one 1-ms subframe in LTE. \\
DCI & \fa{اطلاعات کنترلی فروسو} & Scheduling message on the PDCCH: resources, MCS, HARQ. \\
CQI & \fa{نشانگر کیفیت کانال} & UE report: highest MCS decodable at 10\% BLER. \\
MCS & \fa{طرح مدولاسیون و کدگذاری} & Modulation order + code rate chosen by the scheduler. \\
BSR / PHR & \fa{گزارش وضعیت بافر / حاشیهٔ توان} & Uplink reports: queued data / remaining power. \\
HARQ & \fa{درخواست تکرار خودکار ترکیبی} & Fast retransmission with soft combining. \\
Proportional fair (PF) & \fa{منصفانهٔ تناسبی} & Scheduler maximising $\sum\log R_k$. \\
Multi-user diversity & \fa{تنوع چندکاربره} & Gain from serving users at their fading peaks. \\
\end{keyterms}

% ---------------------------------------------------------------------
\section{Introduction and motivation}
Every millisecond an LTE eNB must answer one question: \emph{who gets which resource blocks,
with which modulation and coding?} The answer determines everything users perceive --
throughput, delay, fairness -- and everything the operator pays for, because spectrum is the
most expensive resource of a mobile network.

The problem is hard because the channel of each user fluctuates (fast fading, mobility,
interference), traffic is bursty, and the objectives conflict: maximising total throughput
favours users close to the base station, while fairness and QoS guarantees require serving
users at the cell edge. LTE gives the eNB a fine-grained tool -- 1\,ms in time, 180\,kHz in
frequency -- and leaves the scheduling \emph{algorithm} to the vendor. 3GPP standardises the
signalling (what the UE reports and how grants are sent), not the decision. This lecture covers
both: the standard mechanisms and the classic algorithms with their mathematical foundations.

% ---------------------------------------------------------------------
\section{Fundamental concepts}

\subsection{The resource grid}
In the time domain an LTE radio frame lasts 10\,ms and contains ten 1\,ms \emph{subframes}, each
made of two 0.5\,ms \emph{slots} of 7 OFDM symbols (normal cyclic prefix) \cite{ts36211}. In the
frequency domain, subcarriers are spaced $\Delta f = 15$\,kHz. A \emph{resource element} is one
subcarrier for one symbol; a \emph{resource block} is 12 subcarriers ($180$\,kHz) for one slot
(84 REs). The scheduler normally assigns \emph{RB pairs} (both slots of a subframe, 168 REs) every
TTI of 1\,ms.

\begin{table}[htbp]
\centering\small
\caption{LTE channel bandwidths and resource blocks \cite{ts36211}.}\label{tab:bw}
\begin{tabular}{@{}lcccccc@{}}
\toprule
Channel bandwidth [MHz] & 1.4 & 3 & 5 & 10 & 15 & 20 \\
Resource blocks $N_{\mathrm{RB}}$ & 6 & 15 & 25 & 50 & 75 & 100 \\
Occupied bandwidth [MHz] & 1.08 & 2.7 & 4.5 & 9.0 & 13.5 & 18.0 \\
\bottomrule
\end{tabular}
\end{table}

\begin{acnfigure}
  \resizebox{\linewidth}{!}{\acnfig{lecture-07/resource-grid}}
  \caption{One downlink RB pair (12 subcarriers $\times$ 14 symbols) with a two-symbol control
  region and two-port cell-specific reference signals.}
  \label{fig:grid}
\end{acnfigure}

\subsection{What the scheduler knows}
\begin{itemize}
  \item \textbf{Channel quality}: CQI reports (wideband or per subband), rank and precoder
        indicators (RI, PMI) in the downlink; sounding reference signals (SRS) in the uplink.
  \item \textbf{Traffic}: downlink buffers are in the eNB; uplink buffers are reported by the UE in
        \emph{buffer status reports} (BSR), per group of logical channels.
  \item \textbf{QoS}: the QCI, guaranteed bit rate (GBR) and priority of each bearer (Lecture~22).
  \item \textbf{UE power}: \emph{power headroom reports} (PHR) tell how much more power the UE could
        use, so the uplink scheduler does not grant more RBs than the UE can power.
\end{itemize}

% ---------------------------------------------------------------------
\section{Detailed technical explanation}

\subsection{Downlink scheduling}
For each scheduled UE the eNB sends a \emph{downlink control information} (DCI) message on the
PDCCH in the control region of the same subframe \cite{ts36213}. A DCI contains the resource
allocation, the MCS index, the HARQ process number, the new-data indicator, the redundancy version
and a power-control command for PUCCH. Three resource allocation types trade signalling size against
flexibility: type~0 (bitmap of RB groups), type~1 (RBs within a subset) and type~2 (a contiguous set,
compact). The UE decodes the PDSCH in the same subframe and returns ACK/NACK four subframes later.

\subsection{Uplink scheduling}
The uplink differs in three ways. (i) The eNB does not see the UE's buffer: the UE sends a
\emph{scheduling request} (SR) on PUCCH, receives a small grant, and then reports its buffer with a
BSR. (ii) Grants are sent on the PDCCH four subframes before the PUSCH transmission (FDD).
(iii) SC-FDMA requires the RBs of one UE to be \emph{contiguous} in frequency to keep the low
peak-to-average power ratio, which limits frequency-selective scheduling. Uplink HARQ is
\emph{synchronous}: retransmissions occur exactly 8\,ms later, acknowledged on the PHICH.

\subsection{CQI and link adaptation}
The UE reports a 4-bit CQI: the highest index whose modulation and code rate would be received
with a transport-block error probability not exceeding 10\% \cite{ts36213}. \Cref{tab:cqi} lists the
standard table. The eNB maps the CQI (after correction, see OLLA below) to an MCS index and a
transport block size.

\begin{table}[htbp]
\centering\small
\caption{LTE 4-bit CQI table (TS 36.213, Table 7.2.3-1). Efficiency is in information bits per
resource element.}\label{tab:cqi}
\rowcolors{2}{acnmist}{white}
\begin{tabular}{@{}cccc|cccc@{}}
\toprule
\rowcolor{white}CQI & Modulation & Code rate$\times1024$ & Efficiency & CQI & Modulation & Code rate$\times1024$ & Efficiency \\
\midrule
1 & QPSK & 78  & 0.1523 & 9  & 16QAM & 616 & 2.4063 \\
2 & QPSK & 120 & 0.2344 & 10 & 64QAM & 466 & 2.7305 \\
3 & QPSK & 193 & 0.3770 & 11 & 64QAM & 567 & 3.3223 \\
4 & QPSK & 308 & 0.6016 & 12 & 64QAM & 666 & 3.9023 \\
5 & QPSK & 449 & 0.8770 & 13 & 64QAM & 772 & 4.5234 \\
6 & QPSK & 602 & 1.1758 & 14 & 64QAM & 873 & 5.1152 \\
7 & 16QAM & 378 & 1.4766 & 15 & 64QAM & 948 & 5.5547 \\
8 & 16QAM & 490 & 1.9141 & 0  & \multicolumn{3}{l}{out of range} \\
\bottomrule
\end{tabular}
\end{table}

Note that the efficiency is simply $Q_m \times (\text{code rate})$: for CQI~15,
$6\times948/1024 = 5.5547$ bits per RE. Wideband CQI describes the whole carrier; subband CQI (per
group of RBs) enables frequency-selective scheduling but costs uplink overhead.

\subsection{HARQ}
LTE uses eight parallel stop-and-wait HARQ processes in FDD (\cref{fig:harq}): while process 2
waits for its acknowledgement, processes 3--8 use the following subframes, so the channel never
idles. With \emph{Chase combining} the retransmission repeats the same coded bits and the receiver
adds the signal energies; with \emph{incremental redundancy} it sends new parity bits, lowering the
effective code rate. Typical operation targets a 10\% first-transmission BLER and a residual BLER
after HARQ of about $10^{-3}$ or below; RLC ARQ (Lecture~6) catches the rest.

\begin{acnfigure}
  \acnfig{lecture-07/harq-processes}
  \caption{Eight HARQ processes in FDD. A NACK for process P2 arrives four subframes after the
  transmission; the retransmission occurs eight subframes after the original.}
  \label{fig:harq}
\end{acnfigure}

\subsection{Scheduling algorithms}
\paragraph{Round robin (RR)} Users take turns regardless of channel quality. Perfectly fair in
\emph{time}, but it ignores the channel and wastes good fading peaks.
\paragraph{Maximum C/I (Max-rate)} In each TTI and RB, serve the user with the highest instantaneous
rate $r_k(t)$. Maximises cell throughput but starves cell-edge users.
\paragraph{Proportional fair (PF)} Serve the user with the largest ratio of instantaneous rate to
its own average throughput \cite{jalali2000hdr}:
\begin{equation}\label{eq:pfmetric}
  k^\ast(t) = \argmax_k \frac{r_k(t)}{\bar R_k(t)}, \qquad
  \bar R_k(t+1) = \Bigl(1-\tfrac{1}{t_\mathrm{c}}\Bigr)\bar R_k(t) + \tfrac{1}{t_\mathrm{c}}\,r_k(t)\,\mathbb{1}\{k = k^\ast(t)\},
\end{equation}
where $t_\mathrm{c}$ is the averaging window in TTIs. Each user is served when its channel is
\emph{good relative to its own average}, so cell-edge users still get their peaks.
\paragraph{QoS-aware scheduling} Real schedulers first satisfy GBR bearers (e.g.\ VoLTE, often with
semi-persistent scheduling), weight metrics by QCI priority and packet delay budget, and use PF-like
rules for the remaining best-effort traffic.

% ---------------------------------------------------------------------
\section{Mathematical models}

\subsection{From CQI to data rate}\label{sec:rate}
For a UE allocated $N_{\mathrm{RB}}$ RB pairs with $N_{\mathrm{RE}}^{\mathrm{data}}$ data REs each and
an efficiency $\eta_{\mathrm{CQI}}$ (bits/RE) from \cref{tab:cqi}, the rate per layer is
\begin{equation}\label{eq:cqirate}
  R = \eta_{\mathrm{CQI}}\,\frac{N_{\mathrm{RB}}\,N_{\mathrm{RE}}^{\mathrm{data}}}{T_{\mathrm{sf}}}.
\end{equation}
\emph{Example:} CQI 9 ($\eta = 2.4063$), 50 RBs, 132 data REs (two control symbols, two CRS ports):
$R = 2.4063\times50\times132/10^{-3} = \qty{15.9}{\mega\bit\per\second}$.

\subsection{Outer-loop link adaptation}\label{sec:olla}
\paragraph{Problem} CQI reports are delayed and biased (estimation errors, interference changes); if
the eNB used them directly the actual BLER would drift from the 10\% target.
\paragraph{Model} The eNB adds an offset $\Delta$ (dB) to the reported SINR before choosing the MCS.
After each ACK it raises $\Delta$ by $\delta_{\mathrm{up}}$; after each NACK it lowers it by
$\delta_{\mathrm{down}}$.
\paragraph{Derivation} In steady state the expected change is zero. If $p$ is the BLER at the
operating point,
\begin{equation}\label{eq:olla}
  (1-p)\,\delta_{\mathrm{up}} - p\,\delta_{\mathrm{down}} = 0
  \quad\Longrightarrow\quad
  p = \frac{\delta_{\mathrm{up}}}{\delta_{\mathrm{up}}+\delta_{\mathrm{down}}}.
\end{equation}
Choosing $\delta_{\mathrm{up}} = \delta_{\mathrm{down}}\,p_\mathrm{t}/(1-p_\mathrm{t})$ therefore drives the
BLER to the target $p_\mathrm{t}$ regardless of the CQI bias. For $p_\mathrm{t} = 0.1$ and
$\delta_{\mathrm{down}} = 0.5$\,dB, $\delta_{\mathrm{up}} = 0.056$\,dB.
\paragraph{Interpretation} OLLA is a stochastic control loop: small steps give an accurate BLER but
slow convergence (about $1/\delta$ steps), large steps react fast but jitter the MCS.

\subsection{HARQ throughput}
If the first transmission fails with probability $p_1$ and a retransmission (after combining) fails
with probability $p_2 < p_1$, the mean number of transmissions per transport block (neglecting a
third attempt's failures) is $\bar N \approx 1 + p_1 + p_1p_2$, and the throughput efficiency is
$1/\bar N$. With $p_1 = 0.1$ and $p_2 = 0.01$: $\bar N = 1.101$, efficiency 90.8\%. A 10\% BLER target
costs only about 9\% of throughput but buys a much higher MCS than a 1\% target would allow -- the
reason LTE operates at 10\%.

\subsection{Why proportional fair is ``proportionally fair''}\label{sec:pf}
\paragraph{Claim} Among all scheduling policies, PF (asymptotically, for long averaging windows)
maximises the utility
\begin{equation}\label{eq:util}
  U = \sum_{k=1}^{K}\log \bar R_k .
\end{equation}
\paragraph{Argument} Consider giving the current TTI to user $k$. To first order it increases
$\bar R_k$ by $\epsilon\,r_k(t)$ (with $\epsilon = 1/t_\mathrm{c}$), changing the utility by
\begin{equation}
  \Delta U \approx \frac{\partial U}{\partial \bar R_k}\,\epsilon\,r_k(t) = \epsilon\,\frac{r_k(t)}{\bar R_k}.
\end{equation}
The greedy choice that maximises the utility increase in every TTI is therefore exactly
$\argmax_k r_k(t)/\bar R_k$, i.e.\ \cref{eq:pfmetric}. Because $U$ is concave, greedy steps converge to
its maximum. A useful property of the optimum is \emph{proportional fairness}: any feasible change of
the throughput vector satisfies $\sum_k \Delta\bar R_k/\bar R_k \le 0$ -- no user can gain a percentage
of throughput without others losing at least the same total percentage. Max C/I maximises $\sum \bar R_k$
instead, and RR equalises time shares.

\subsection{Jain's fairness index}
To compare schedulers we use \cite{jain1984fairness}
\begin{equation}\label{eq:jain}
  J = \frac{\bigl(\sum_{k=1}^{K}\bar R_k\bigr)^2}{K\sum_{k=1}^{K}\bar R_k^2}, \qquad \frac{1}{K}\le J\le 1,
\end{equation}
which equals 1 when all users get the same throughput and $1/K$ when one user gets everything.

\subsection{Simulation: multi-user diversity}\label{sec:sim}
\paragraph{Set-up} \texttt{scripts/scheduler\_sim.py} simulates one 10\,MHz carrier (9\,MHz occupied)
given entirely to one user per TTI; four users with mean SNRs of 20, 12, 5 and 0\,dB; independent
Rayleigh block fading per TTI; efficiency $\min(0.75\log_2(1+\SNR),5.55)$ bit/s/Hz (an attenuated
Shannon mapping capped at CQI~15); 20\,000 TTIs; $t_\mathrm{c} = 100$; fixed random seed.
\paragraph{Results} (\cref{fig:sched}, \cref{tab:sim})
\begin{table}[htbp]
\centering\small
\caption{Simulated per-user and cell throughput [Mbit/s] and Jain's index.}\label{tab:sim}
\begin{tabular}{@{}lrrrrrr@{}}
\toprule
Scheduler & UE1 (20\,dB) & UE2 (12\,dB) & UE3 (5\,dB) & UE4 (0\,dB) & Cell & $J$ \\
\midrule
Round robin  & 9.74  & 5.83 & 2.88 & 1.43 & 19.88 & 0.710 \\
Prop.\ fair  & 12.97 & 8.34 & 4.60 & 2.50 & 28.40 & 0.761 \\
Max C/I      & 35.73 & 4.18 & 0.21 & 0.01 & 40.13 & 0.311 \\
\bottomrule
\end{tabular}
\end{table}
\paragraph{Interpretation} PF gives \emph{every} user more throughput than RR: each user still gets
about a quarter of the TTIs, but they are its best TTIs. This is \emph{multi-user diversity}, a gain
that grows with the number of users and the fading depth. Max C/I raises the cell total by another 41\%
but essentially switches off UE3 and UE4. The model ignores frequency selectivity (subband
scheduling adds further diversity), CQI delay and errors, and traffic dynamics (all users always have
data -- the ``full buffer'' assumption).

\begin{acnfigure}
  \acnfig{lecture-07/scheduler-results}
  \caption{Per-user throughput under round robin, proportional fair and maximum C/I scheduling
  (simulation of \cref{sec:sim}).}
  \label{fig:sched}
\end{acnfigure}

% ---------------------------------------------------------------------
\section{Network architecture and diagrams}
Scheduling is entirely a function of the eNB MAC; the only network-level inputs are the QoS
parameters of each bearer, set by the core (QCI, ARP, GBR, MBR, AMBR) through S1AP when bearers are
established (Lectures~9 and~22). Between eNBs, X2 \emph{load information} messages allow inter-cell
interference coordination: neighbours can agree to avoid scheduling cell-edge users on the same RBs.
\Cref{fig:trilemma} places the algorithms in the space of objectives.

\begin{acnfigure}
  \acnfig{lecture-07/trilemma}
  \caption{The scheduling trilemma. No algorithm maximises throughput, fairness and QoS at once;
  practical schedulers combine a QoS stage with a PF-like stage.}
  \label{fig:trilemma}
\end{acnfigure}

% ---------------------------------------------------------------------
\section{Step-by-step procedures}
\begin{procedurebox}{One TTI of downlink scheduling}\label{proc:dl}
\begin{steps}
  \item Collect inputs: latest CQI/RI/PMI per UE, HARQ feedback received for subframe $n-4$, buffer
        sizes per bearer, QoS parameters.
  \item Reserve resources for pending HARQ retransmissions (they have priority to keep delay low).
  \item Serve GBR bearers that are below their guaranteed rate (e.g.\ VoLTE), using the MCS from
        corrected CQI (OLLA offset).
  \item Rank the remaining backlogged UEs with the PF metric per RB group and allocate RBs.
  \item For each scheduled UE choose MCS and transport block size; build the MAC PDUs (multiplexing
        logical channels by priority).
  \item Encode the DCIs into PDCCH candidates (if the control region is full, defer UEs to the next TTI)
        and transmit PDCCH and PDSCH.
  \item Update PF averages and OLLA offsets when ACK/NACK arrive four subframes later.
\end{steps}
\end{procedurebox}

\begin{procedurebox}{Uplink access for a UE with new data}\label{proc:ul}
\begin{steps}
  \item New data arrives in an empty buffer; the UE triggers a BSR. Having no grant, it sends a
        scheduling request on its next PUCCH SR opportunity (or random access if none is configured).
  \item The eNB sends a small grant (DCI format 0); four subframes later the UE transmits the BSR (and
        PHR) on PUSCH.
  \item The eNB schedules the UE according to the reported buffer, its uplink SINR (from SRS/DMRS) and
        its power headroom, with contiguous RBs.
  \item Each PUSCH transmission is acknowledged on PHICH after four subframes; non-adaptive
        retransmissions follow eight subframes after the original.
\end{steps}
\end{procedurebox}

% ---------------------------------------------------------------------
\section{Worked examples}

\begin{examplebox}{Counting resources}\label{ex:count}
\textbf{Problem.} How many REs does a 15\,MHz LTE carrier offer per second, and how many of them are
PDSCH REs with a three-symbol control region and four CRS ports?

\textbf{Solution.} $75\times12\times14\times1000 = 12.6\times10^6$ RE/s. Data REs per RB pair: $168 - 36 -
16 = 116$ (four CRS ports use 24 REs, 8 of which lie in the three control symbols). Total
$75\times116\times1000 = 8.7\times10^6$ RE/s, i.e.\ 69\%.
\end{examplebox}

\begin{examplebox}{OLLA step}\label{ex:olla}
\textbf{Problem.} An operator wants a 20\% BLER target for a delay-tolerant service and uses
$\delta_{\mathrm{down}} = 1$\,dB. Find $\delta_{\mathrm{up}}$ and the offset drift if the actual BLER were
10\%.

\textbf{Solution.} $\delta_{\mathrm{up}} = 1\times0.2/0.8 = 0.25$\,dB. At 10\% BLER the expected change per
transmission is $0.9\times0.25 - 0.1\times1 = +0.125$\,dB: the offset rises (more aggressive MCS) until
the BLER reaches 20\%.
\end{examplebox}

\begin{examplebox}{PF decision}\label{ex:pf}
\textbf{Problem.} Two UEs: UE~A has average 10\,Mbit/s and instantaneous 18\,Mbit/s; UE~B has average
2\,Mbit/s and instantaneous 5\,Mbit/s. Whom do Max C/I and PF serve?

\textbf{Solution.} Max C/I: UE~A (18 > 5). PF: $18/10 = 1.8$ versus $5/2 = 2.5$ $\Rightarrow$ UE~B, which
is at a fading peak relative to its own average.
\end{examplebox}

% ---------------------------------------------------------------------
\section{Technology comparison}
\begin{table}[htbp]
\centering\small
\caption{Scheduling and link adaptation across generations.}\label{tab:gens}
\rowcolors{2}{acnmist}{white}
\begin{tabularx}{\linewidth}{@{}l Y Y Y@{}}
\toprule
\rowcolor{white} & HSDPA (3.5G) & LTE (4G) & NR (5G) \\
\midrule
Scheduling unit   & 2\,ms TTI, codes            & 1\,ms TTI, RB pairs (180\,kHz)  & slot (1--0.125\,ms) or mini-slot, RBs of $12\Delta f$ \\
Frequency domain  & none (CDMA)                 & per RB group (OFDMA)            & per RB, bandwidth parts \\
Channel feedback  & CQI                         & CQI, PMI, RI (subband)          & CSI-RS based, beam reports \\
HARQ              & NodeB, 6 processes          & 8 processes, 8\,ms RTT (FDD)   & up to 16 processes, flexible timing \\
Uplink            & E-DCH (Rel-6)               & SC-FDMA, contiguous RBs         & CP-OFDM or DFT-s-OFDM, configured grants \\
\bottomrule
\end{tabularx}
\end{table}

\paragraph{Design trade-offs} Finer granularity enables better adaptation but costs control overhead
(more DCIs, more CQI reports); frequent feedback improves accuracy but uses uplink capacity; longer PF
windows favour throughput, shorter ones favour short-term fairness and delay. 5G extends the same
principles with flexible numerology (Lecture~14).

% ---------------------------------------------------------------------
\section{Key takeaways}
\begin{takeawaybox}
\begin{itemize}
  \item Frame 10\,ms $\to$ subframe 1\,ms $\to$ slot 0.5\,ms (7 symbols); RB $=$ 12 subcarriers $\times$ slot;
        $N_{\mathrm{RB}} = 6\dots100$ for 1.4--20\,MHz.
  \item 3GPP standardises signalling (DCI, CQI, BSR, PHR, HARQ timing), not the scheduling algorithm.
  \item CQI $\to$ efficiency $Q_m\times$ code rate; OLLA drives the BLER to $\delta_{\mathrm{up}}/(\delta_{\mathrm{up}}+\delta_{\mathrm{down}})$.
  \item Eight HARQ processes keep the channel busy; a 10\% first-transmission BLER costs about 9\%.
  \item Max C/I maximises $\sum \bar R_k$, PF maximises $\sum\log\bar R_k$, RR equalises time.
  \item Multi-user diversity: PF can beat RR for \emph{every} user in fading channels.
  \item Practical schedulers: HARQ retransmissions first, then GBR/QoS, then PF for best effort.
\end{itemize}
\end{takeawaybox}

% ---------------------------------------------------------------------
\section{Exercises}

\begin{exercisebox}{1}{Standard versus vendor}\label{exr:std}
Which of the following are standardised by 3GPP and which are vendor implementation? CQI table,
proportional fair metric, DCI format, OLLA step sizes, HARQ timing in FDD, priority of GBR bearers.
\end{exercisebox}

\begin{exercisebox}{1}{Uplink contiguity}\label{exr:contig}
Explain why SC-FDMA requires contiguous RBs per UE and what this means for frequency-selective
uplink scheduling.
\end{exercisebox}

\begin{exercisebox}{2}{Rate from CQI}\label{exr:cqi}
A UE reports CQI 12 with rank 2 on a 20\,MHz carrier (one control symbol, two CRS ports). It receives 40
RBs. Compute its rate with \cref{eq:cqirate}.
\end{exercisebox}

\begin{exercisebox}{2}{Jain's index}\label{exr:jain}
Verify the Jain indices of \cref{tab:sim} for RR and Max C/I from the per-user values.
\end{exercisebox}

\begin{exercisebox}{3}{Two-user PF analysis}\label{exr:pf2}
Two users have i.i.d.\ channels with the same distribution but user~1's rate is always twice user~2's
($r_1 = 2r_2$ in every TTI). Show that PF gives each user half of the TTIs, and compute the resulting
throughputs relative to RR. What changes if the channels are independent?
\end{exercisebox}

\begin{exercisebox}{3}{Simulation study}\label{exr:sim}
Modify \texttt{scheduler\_sim.py} to use 8 users with mean SNRs uniformly spaced between 0 and 21\,dB and
plot cell throughput and Jain's index versus $t_\mathrm{c}\in\{10,100,1000\}$ for PF. Explain the trend.
\end{exercisebox}

\subsection*{Solutions}
\addcontentsline{toc}{subsection}{Solutions}

\begin{solutionbox}{\ref{exr:std}}
Standardised: CQI table, DCI formats, HARQ timing in FDD (and the meaning of BSR/PHR). Vendor
implementation: the PF metric (or any algorithm), OLLA step sizes, the prioritisation policy among
bearers (constrained by the QoS parameters but not dictated in detail).
\end{solutionbox}

\begin{solutionbox}{\ref{exr:contig}}
SC-FDMA is DFT-spread OFDM: mapping the DFT output to contiguous subcarriers produces a
single-carrier-like signal with low peak-to-average power ratio, allowing efficient power amplifiers
in the handset. Non-contiguous mapping raises the PAPR. Hence the scheduler can choose \emph{where}
a UE's block lies (frequency-selective gain) but not pick scattered good RBs; Release~10 relaxed this
with clustered allocation of two clusters.
\end{solutionbox}

\begin{solutionbox}{\ref{exr:cqi}}
Data REs per RB pair: $168 - 12 - (16-4) = 144$. $R = 2\times3.9023\times40\times144/10^{-3} =
\qty{45.0}{\mega\bit\per\second}$.
\end{solutionbox}

\begin{solutionbox}{\ref{exr:jain}}
RR: sum $=19.88$, sum of squares $=94.87+33.99+8.29+2.04 = 139.2$; $J = 395.2/(4\times139.2) = 0.710$.
Max C/I: sum $= 40.13$, squares $\approx 1276.6+17.5+0.04 = 1294.2$; $J = 1610.4/(4\times1294.2) = 0.311$.
\end{solutionbox}

\begin{solutionbox}{\ref{exr:pf2}}
With $r_1 = 2r_2$ always, $r_1/\bar R_1$ versus $r_2/\bar R_2$ compares $2r_2/\bar R_1$ with $r_2/\bar R_2$;
equilibrium requires both to be chosen, which happens when $\bar R_1 = 2\bar R_2$, i.e.\ when each user
gets half the TTIs -- exactly RR, because the channels are perfectly correlated and there is no
multi-user diversity to exploit. With independent channels PF picks each user at its relative peaks,
and both throughputs exceed RR, as in \cref{tab:sim}.
\end{solutionbox}

\begin{solutionbox}{\ref{exr:sim}}
With mean SNRs $0,3,\dots,21$\,dB and 40\,000 TTIs the script gives: $t_\mathrm{c}=10$: cell 29.8\,Mbit/s,
$J=0.831$; $t_\mathrm{c}=100$: 32.8\,Mbit/s, $J = 0.838$; $t_\mathrm{c}=1000$: 33.2\,Mbit/s, $J=0.839$ (RR:
21.8\,Mbit/s, $J=0.775$). A short window reacts to recent service and sometimes serves a user who is not
at a relative peak; longer windows let PF wait for genuine peaks, so throughput rises and saturates once
the window is much longer than the fading (here independent from TTI to TTI). The price of long windows
in practice is short-term unfairness and larger delay jitter when channels are time-correlated.
\end{solutionbox}

% ---------------------------------------------------------------------
\section{Further reading}
The physical channels and grid are specified in \cite{ts36211}, CQI, DCI and HARQ procedures in
\cite{ts36213} and MAC procedures (BSR, PHR, SR) in \cite{ts36321}. The proportional fair scheduler
originated in CDMA-HDR \cite{jalali2000hdr}; its multi-user diversity interpretation is developed in
Tse and Viswanath \cite{tse2005fundamentals}. Jain's index is from \cite{jain1984fairness}.
Sesia et al.\ \cite{sesia2011lte} cover LTE scheduling in depth.

\printbibliography[heading=subbibliography,title={\ACNstrRefs}]
\end{document}
